\section{问题三模型的建立与求解}

\subsection{问题三求解思路}
问题三在问题二的基础上增加了每天 0:00、6:00、12:00、18:00 发布的未来 24 小时整点光伏预报（附件 3），
并允许在 6:00、12:00、18:00 依据新预报调整当天尚未执行的购电计划：少购部分按交易时刻电价的 50\% 结算，
多购部分按 150\% 结算，供电缺口仍按 5 倍电价紧急购电。因此，本问的关键并不在于“用最新预报直接替换原计划”，
而在于回答\textbf{后续时刻的预报是否值得用于制定调整购电策略}。

为此，本文首先对附件 3 的预报质量进行诊断，判断后续预报相对早前预报是否带来新的信息；
在此基础上建立“预报标定—当日残差修正—双向滚动调整—末端备用”的两阶段调度模型，
通过策略变体对比拆解“预报质量”与“调整机制”各自的贡献，并定量回答上述问题。
问题三的求解流程可概括为“日前计划—日内滚动调整—实际运行结算”三个阶段，如图~\ref{fig:q3-flow} 所示。

\begin{figure}[!htbp]
\centering
\footnotesize
\begin{tikzpicture}[
  >=Stealth,
  node distance=7mm,
  plan/.style={rectangle, draw=black!70, fill=orange!12, align=left, text width=140mm,
               minimum height=10mm, inner sep=3pt},
  stage/.style={rectangle, draw=black!70, fill=blue!6, align=left, text width=140mm,
               minimum height=10mm, inner sep=3pt},
  io/.style={rectangle, rounded corners=2mm, draw=black!70, fill=gray!12, align=center,
               text width=140mm, minimum height=8mm, inner sep=3pt},
  line/.style={draw, -Latex, thick}
]
  \node[plan] (p1) {\textbf{0:00　制定日前计划}：将 MOS 标定后的光伏估计与上周同星期实际值（w1）融合，
                    并按预测余量 $\gamma=1.02$、$\delta=0.98$ 构造保守净负荷；\\
                    求解 $\min\sum_{t} c_t p^{plan}_t\Delta t$，
                    满足功率平衡、储能容量与功率约束、日循环约束 $S^{plan}_{144}=S^{real}_{0}$};
  \node[stage] (a1) [below=of p1] {\textbf{6:00 / 12:00 / 18:00　双向滚动调整}：接收新发布的整点预报
                    $\rightarrow$ MOS 标定 $+$ 当日已实现残差修正（$\lambda_{pv}=0.1$、$\lambda_L=0.3$）；\\
                    求解调整 LP：$\min\sum_{t\ge a}\left[c_t p^{adj}_t+0.5c_t(u^{+}_t+u^{-}_t)+5c_t e_t\right]\Delta t$，
                    末端最低备用 $S^{plan}_{144}\ge 4000$\,kWh};
  \node[stage] (r1) [below=of a1] {\textbf{全天 144 个 $10\,$min　日内实际运行}：实际负载与光伏
                    $\rightarrow$ 储能实时纠偏 $\rightarrow$ 缺口按 $5c_t$ 紧急购电、富余弃光；\\
                    按 $C(t)=c_t p^{adj}_t+0.5c_t\left|p^{adj}_t-p^{plan}_t\right|+5c_t e_t$ 结算当段购电费用};
  \node[io] (o1) [below=of r1] {全年购电总费用与四个指定日期的计划、调整结果
                    $\rightarrow$ 与问题二策略及完美预测下界的对比};

  \draw[line] (p1.south) -- node[right, font=\scriptsize] {计划购电量 $p^{plan}_t$} (a1.north);
  \draw[line] (a1.south) -- node[right, font=\scriptsize] {调整后购电量 $p^{adj}_t$} (r1.north);
  \draw[line] (r1.south) -- (o1.north);
\end{tikzpicture}
\caption{问题三“计划—调整—结算”两阶段滚动调度流程}
\label{fig:q3-flow}
\end{figure}

\subsection{问题三模型建立}

\subsubsection{附件 3 预报质量诊断}
为判断后续预报的信息价值，将附件 3 中全年 365 天、每天 4 次发布的整点预报与附件 2 的实际光伏出力
按同一口径逐时对照（预报“第 $k$ 小时”对应发布时刻后第 $k$ 个整点，实际值取该小时内 6 个 $10\,$min 区间的平均值），
统计各目标时刻的平均绝对误差（MAE）\cite{fpp3} 第 5.4 节 与系统性偏差，结果如图~\ref{fig:q3-forecast} 所示。

\begin{figure}[!htbp]
\centering
\includegraphics[width=0.485\textwidth]{Q3_mae.pdf}\hfill
\includegraphics[width=0.485\textwidth]{Q3_bias.pdf}
\caption{附件 3 预报质量诊断：（左）不同发布时刻的预报 MAE 随目标时刻的变化；
（右）预报偏差（预报值$-$实际值）随目标时刻的变化}
\label{fig:q3-forecast}
\end{figure}

由图~\ref{fig:q3-forecast}（左）可知，\textbf{同一目标时刻的预报精度几乎不随发布时间改善}：
目标时刻为 16:00 时，0:00、6:00、12:00 发布的预报 MAE 分别为 1000.2、1043.6、1066.5\,kW，
三者基本相当；而在 14:00 之后，上周同星期实际值（w1）的 MAE 反而显著更低（16:00 时仅 214.7\,kW）。
12:00 发布的预报在目标时刻 12:00 处的 MAE 高达 1374.1\,kW，明显劣于当日 0:00 的预报（640.9\,kW），
说明\textbf{“更新”的预报并不必然优于早前预报}。

由图~\ref{fig:q3-forecast}（右）可知，附件 3 的预报存在明显的\textbf{系统性偏差}：
上午时段预报偏低（8:00 约 $-1050$\,kW，10:00 约 $-580\sim-700$\,kW），
午后时段预报偏高（14:00 约 $+567\sim+790$\,kW，15:00 约 $+836\sim+988$\,kW，16:00 约 $+988\sim+1067$\,kW），
且不同发布时刻的偏差方向高度一致，表明这是预报本身的系统性偏差而非某次发布的偶然误差。
此外，预报峰值出现的时刻（13:00）与实际峰值时刻（12:00）存在错位。

进一步计算可知，尽管后续预报的信息增量有限，预报值与实际光伏出力仍具有较高的相关性（逐时相关系数 0.78$\sim$0.91），
其系统性偏差可通过历史数据标定加以修正；同时，当日已实现的预报偏差对后续时段具有持续性：
按白天整点块均值统计，当日上午实际出力与午后实际出力的相关系数为 0.914，
w1 残差的日内自相关系数光伏为 0.464、负载为 0.592。
这说明\textbf{后续预报的正确用法是“经标定的输入与当日已实现残差的修正”，而不是当作真值直接使用}。

\subsubsection{光伏出力估计模型}
考虑到预报存在系统性偏差，本文采用模型输出统计（MOS）方法\cite{glahn1972}对光伏出力进行标定，
并将当日已实现残差用于对后续时段的修正。

\textbf{（1）0:00 计划用估计值。} 逐小时建立线性回归\cite{fpp3} 第 7 章，将当日 0:00 官方预报与上周同星期实际值（w1）融合：
\begin{equation}
\widehat{PV}(d,h)=a_h+b_h\,\widehat{PV}^{off}_{0:00}(d,h)+c_h\,\widehat{PV}^{w1}(d,h)
\label{eq:q3-mos}
\end{equation}
其中回归系数 $a_h,b_h,c_h$ 由最近 30 天的（官方预报，w1，实际值）三元组按最小二乘拟合，
并逐日滚动更新以保证因果性；历史数据不足 5 天时退化为两者等权（0.5/0.5）融合。
例如 14:00 的估计式为 $\widehat{PV}\approx 232+0.34\,\widehat{PV}^{off}_{0:00}+0.58\,\widehat{PV}^{w1}$，
其 MAE 为 373\,kW，优于单独使用 w1（456.8\,kW）和使用原始预报（793.4\,kW）。
逐小时估计值再按上周同日的日内形状系数分解到 144 个 $10\,$min 时段。

\textbf{（2）调整时刻的更新估计值。} 在 6:00、12:00、18:00 三个调整时刻 $\tau$，
将 $\tau$ 时刻发布的官方预报代入式~\ref{eq:q3-mos} 得到基准估计，并叠加当天已实现残差：
\begin{equation}
\begin{aligned}
\widehat{PV}_\tau(h)&=\widehat{PV}(h)+\lambda_{pv}\overline{\Big(PV^{real}-\widehat{PV}\Big)}_{h'<\tau}\\
\widehat{L}_\tau(h)&=\widehat{L}(h)+\lambda_L\overline{\Big(L^{real}-\widehat{L}\Big)}_{h'<\tau}
\end{aligned}
\label{eq:q3-resid}
\end{equation}
其中 $\overline{(\cdot)}_{h'<\tau}$ 表示对 $\tau$ 时刻之前已实现的整点时段取平均，
$\lambda_{pv}$、$\lambda_L$ 为残差更新权重。由于残差的日内自相关系数为正但不强，
$\lambda$ 宜取中小值；本文以全年总费用为目标在网格上对二者寻优，
最终取 $\lambda_{pv}=0.1$、$\lambda_L=0.3$。

\subsubsection{计划与调整两阶段优化模型}
\textbf{（1）费用结算口径。} 按题意，第 $t$ 个 $10\,$min 区间的购电费用为
\begin{equation}
C(t)=c_t\,p^{adj}_t+0.5\,c_t\left|p^{adj}_t-p^{plan}_t\right|+5\,c_t\,e_t
\label{eq:q3-cost}
\end{equation}
式~\ref{eq:q3-cost} 中三项依次为按调整购电量实付的电费、调整购电量相对计划量的违约与溢价费用、
以及供电缺口部分的紧急购电费用；该分段线性费用精确等价于：计划内部分按 $c_t$ 结算、
超出计划部分按 $1.5c_t$ 结算、计划内未购足部分按 $0.5c_t$ 结算。为便于线性规划求解，
引入非负变量 $u^{+}_t$、$u^{-}_t$ 表示相对计划的上调与下调购电量：
\begin{equation}
p^{adj}_t-p^{plan}_t=u^{+}_t-u^{-}_t,\qquad u^{+}_t\ge 0,\quad u^{-}_t\ge 0
\label{eq:q3-split}
\end{equation}

\textbf{（2）0:00 计划模型。} 其结构与问题二一致：以负载预测上调 2\%、光伏预测下调 2\% 的预测余量
（$\gamma=1.02$、$\delta=0.98$）构造保守净负荷，以当日计划购电费用最小为目标，
满足功率平衡、储能容量与充放电功率约束以及日循环约束 $S_{144}=S_0$，
得到计划购电量 $p^{plan}_t$ 与计划储能轨迹。

\textbf{（3）调整模型。} 在调整时刻 $\tau\in\{6,12,18\}$，以当日 0:00 计划为基准，
对尚未执行的时段 $t\ge a$ 重新优化，其中 $a$ 为该调整时刻对应的时段编号：
\begin{equation}
\min\ \sum_{t=a}^{144}\Big[c_t\,p^{adj}_t+0.5\,c_t\,u^{+}_t+0.5\,c_t\,u^{-}_t+5\,c_t\,e_t\Big]\Delta t
\label{eq:q3-adjlp}
\end{equation}
约束条件包括：更新估计 $\widehat{L}_\tau$、$\widehat{PV}_\tau$ 下的功率平衡关系，
以该时刻实际储电量为初值的储能状态转移方程与容量、功率约束，式~\ref{eq:q3-split} 的调整量分解关系，
以及紧急购电松弛变量 $e_t\ge 0$。引入 $e_t$ 作为“安全阀”，一方面保证任意更新情形下调整模型可行，
另一方面使模型能在“按 $1.5c_t$ 上调购电”与“按 $5c_t$ 紧急购电”之间自动比价。
此外，为避免模型在傍晚短视放空储能，在调整模型中设置\textbf{日末最低备用约束}：
\begin{equation}
S^{plan}_{144}\ge S^{end}_{\min}=4000\ \text{kWh}
\label{eq:q3-end}
\end{equation}
该约束仅作用于调整计划的规划口径；日内实际偏差下，储能仍可放电至运行下限（1200\,kWh）。

\textbf{（4）日内实际运行与结算。} 日内仍采用与问题二相同的储能实时纠偏机制
（缺额时依次减少计划充电、增加储能放电、按 5 倍电价紧急购电；
富余时依次减少计划放电、增加储能充电、弃置剩余电量），并按式~\ref{eq:q3-cost} 对调整后的实际购电量结算。

\subsection{问题三模型求解与分析}

\subsubsection{参数寻优}
\label{sec:q3-param}
调整机制的存在使预测偏差的代价由问题二中的 5 倍紧急购电下降为 0.5 倍调整费用，
因此预测余量参数不再沿用问题二的取值。以全年总费用为指标重新寻优，结果表明：
负载余量系数的最优值由 1.05 降至 $\gamma=1.02$，光伏余量系数保持 $\delta=0.98$。

对当日残差更新权重 $\lambda$ 的寻优结果如表~\ref{tab:q3-lambda} 所示。
不利用当日信息（$\lambda=0$）时全年总费用为 1376.3 万元，
在 $\lambda_{pv}=0.1$、$\lambda_L=0.3$ 处取得最小值 1368.8 万元，
说明当日已实现的负载与光伏偏差确实提供了有效的信息增量；
而当 $\lambda$ 取值偏大（如 0.3/0.4 以上）时，调整模型会追逐残差中的随机噪声、支付多余的调整费用，
总费用反而回升。最优解附近曲面平坦，说明结果对 $\lambda$ 的具体取值不敏感。

\begin{table}[!htbp]
\caption{当日残差更新权重的寻优结果（全年总费用，单位：万元）}\label{tab:q3-lambda}
\centering
\small
\begin{tabular}{lccccccc}
\toprule
$(\lambda_{pv},\lambda_L)$ & $(0,0)$ & $(0.1,0.2)$ & $(0.1,0.3)$ & $(0.15,0.2)$ & $(0.3,0.4)$ & $(0.45,0.6)$ & $(0.6,0.8)$ \\
\midrule
全年总费用 & 1376.3 & 1369.5 & \textbf{1368.8} & 1370.6 & 1377.2 & 1393.0 & 1414.4 \\
\bottomrule
\end{tabular}
\end{table}

末端最低备用约束的取值同样经变体实验确定：完全自由末端虽比强制当日闭环略省费用
（1382.6 万元对 1384.3 万元，仅相差 0.12\%），但会使储能每天傍晚放空至运行下限附近、次日再低价回充，
几乎不保留备用容量，全年紧急购电天数明显增加；设置 $S^{end}_{\min}=4000$\,kWh
（约为可用容量的 42\%）后，既抑制了末端短视行为，又为次日保留了基本备用。
该约束在 2000$\sim$7000\,kWh 范围内全年总费用变化平缓，说明结果对具体取值不敏感。

\subsubsection{策略变体对比}
为定量回答“是否需要引入其他时刻的预报”，本文设计了一组策略变体进行全年滚动仿真，
从“预报质量”与“调整机制”两个维度逐步改进，结果如表~\ref{tab:q3-variant} 和图~\ref{fig:q3-scheme} 所示。

\begin{table}[!htbp]
\caption{不同购电策略的全年总费用对比（报告期 2025 年 2 月 1 日--12 月 31 日，共 334 天）}\label{tab:q3-variant}
\centering
\small
\setlength{\tabcolsep}{4pt}
\begin{tabular}{lllll}
\toprule
策略 & 说明 & 总费用/万元 & 费用构成/万元 & 相对问题二 \\
\midrule
V0 & 问题二策略（w1 光伏预测、无调整） & 1454.7 & 实付 1367.6 $+$ 紧急 87.0 & --- \\
V1 & 原始 0:00 预报、无调整 & 1809.5 & 实付 1335.7 $+$ 紧急 473.8 & $+$24.4\% \\
V2 & 原始预报 $+$ 双向调整 & 1586.0 & 实付 1458.3 $+$ 调整 110.1 $+$ 紧急 17.7 & $+$9.0\% \\
V3 & 预报标定、无调整 & 1410.7 & 实付 1315.0 $+$ 紧急 95.6 & $-$3.0\% \\
V4D & 标定 $+$ 仅下调调整 & 1407.0 & 实付 1301.4 $+$ 违约 17.5 $+$ 紧急 88.1 & $-$3.3\% \\
V4 & 标定 $+$ 双向调整（强制当日闭环） & 1384.3 & 实付 1312.8 $+$ 调整 33.7 $+$ 紧急 37.8 & $-$4.8\% \\
V4F & 标定 $+$ 双向调整 $+$ 自由末端 & 1382.6 & 实付 1303.6 $+$ 调整 36.3 $+$ 紧急 42.6 & $-$5.0\% \\
\textbf{V4FR} & \textbf{标定 $+$ 双向调整 $+$ 末端备用 $+$ $\lambda$ 寻优} & \textbf{1368.8} & \textbf{实付 1306.2 $+$ 调整 32.1 $+$ 紧急 30.4} & \textbf{$-$5.9\%} \\
--- & 完美预测（理论下界） & 1224.5 & --- & $-$15.8\% \\
\bottomrule
\end{tabular}
\end{table}

\begin{figure}[!htbp]
\centering
\includegraphics[width=0.82\textwidth]{Q3_scheme.pdf}
\caption{各策略变体的全年总费用对比（虚线为完美预测下的理论下界 1224.5 万元）}
\label{fig:q3-scheme}
\end{figure}

由表~\ref{tab:q3-variant} 可知：直接使用原始预报替代 w1 预测会使总费用由 1454.7 万元上升至 1809.5 万元
（$+$24.4\%），说明\textbf{原始预报的偏差远大于 w1 预测}；
但仅依靠调整机制（V2）即可挽回其中 63\% 的损失，
说明\textbf{调整机制本身具有可观的经济价值}——它把 5 倍电价的紧急购电转化为 1.5 倍溢价、
把过配导致的全部损失削减为 50\% 的违约费用。进一步引入预报标定（V3）后，
总费用降至 1410.7 万元，说明\textbf{标定是发挥预报价值的前提}；
在标定基础上再引入双向调整、末端备用与残差更新权重寻优后（V4FR），
总费用最终降至 1368.8 万元，较问题二策略降低 5.9\%。

对比 V4 与 V4F 可知，末端处理方式本身的影响有限（相差 1.7 万元），
但完全自由末端会使储能长期贴近运行下限、削弱备用能力，故本文采用带末端最低备用约束的 V4FR；
对比 V4D 与 V4 可知，仅允许下调购电量会因无法应对光伏低估而多支付约 50 万元紧急购电费用
（88.1 万元对 37.8 万元），说明\textbf{双向调整是必要的}。

\subsubsection{求解结果与分析}
本文最终策略的全年运行指标如表~\ref{tab:q3-annual} 所示。报告期内微网实际购电总费用为 1368.75 万元，
较问题二策略（1454.65 万元）降低 85.90 万元（$-$5.9\%）；
其中计划购电费用为 1286.32 万元，按实际调整购电量结算的实付费用为 1306.21 万元，
调整违约与溢价费用为 32.14 万元（上调购电量 61.60 万\,kWh、下调 38.23 万\,kWh），
紧急购电费用为 30.40 万元，较问题二策略的 87.02 万元下降 65.1\%。
这表明日内滚动调整有效吸收了预测偏差，把原本需以 5 倍电价兜底的电力缺口转化为代价较低的调整购电；
全年紧急购电量也由问题二策略的 14.51 万\,kWh 降至 5.59 万\,kWh。

\begin{table}[!htbp]
\caption{问题三最终策略的年度运行指标（报告期 2025 年 2 月 1 日--12 月 31 日）}\label{tab:q3-annual}
\centering
\small
\setlength{\tabcolsep}{5pt}
\begin{tabular}{lrr}
\toprule
指标 & 数值 & 占比 \\
\midrule
计划购电费用（0:00 计划口径） & 12863249.43 元 & --- \\
按调整购电量实付费用 & 13062110.51 元 & 100.0\% \\
调整违约与溢价费用 & 321415.34 元 & 2.5\% \\
紧急购电费用（5 倍电价） & 303964.84 元 & 2.3\% \\
\midrule
全年购电总费用 & \textbf{13687490.69 元（1368.75 万元）} & --- \\
\bottomrule
\end{tabular}
\end{table}

四个指定日期的计划与调整结果如表~\ref{tab:q3-days} 所示。四个日期的调整均集中在
12:00--12:10 与 18:00--18:10 两个时段：3 月 20 日 12:00 由计划 629.79\,kWh 上调至 839.81\,kWh；
6 月 21 日 18:00 由 444.99\,kWh 微调至 479.30\,kWh；9 月 23 日 12:00 由 529.25\,kWh 上调至 852.32\,kWh；
12 月 21 日 12:00 由 996.21\,kWh 上调至 1256.55\,kWh、18:00 由 731.57\,kWh 微调至 736.34\,kWh；
其余四个指定时段（10:00、14:00、16:00、20:00）的购电量与计划一致。
四个日期 0:00 与 24:00 的储电量均处于 4100$\sim$5200\,kWh 的健康备用区间，
说明末端备用约束有效避免了储能被放空。

\begin{table}[!htbp]
\caption{四个指定日期的计划购电量、调整后购电量、购电费用与紧急购电情况}\label{tab:q3-days}
\centering
\small
\setlength{\tabcolsep}{5pt}
\begin{tabular}{lcccccc}
\toprule
日期 & 计划全天/kWh & 调整全天/kWh & 计划购电费/元 & 调整后结算/元 & 紧急购电 \\
\midrule
2025-03-20 & 70931.49 & 71641.61 & 43191.79 & 43970.52 & 无 \\
2025-06-21 & 35382.09 & 35651.35 & 20188.83 & 21095.57 & 无 \\
2025-09-23 & 67776.12 & 69583.76 & 42138.53 & 44294.40 & 20:30--20:40，6.38\,kWh \\
2025-12-21 & 96674.20 & 98469.77 & 61965.50 & 63563.48 & 8:00--8:10，29.55\,kWh \\
\bottomrule
\end{tabular}
\end{table}

\subsubsection{是否需要引入其他时刻的预报}
综合上述诊断与对比结果，对本题所提问题的回答是：\textbf{需要，但必须以正确的方式引入}。

\textbf{（1）后续预报本身的信息增量有限。} 同一目标时刻的预报精度并不随发布时间改善，
其修订量与当日已实现的偏差也几乎不相关，因此\textbf{不宜用最新预报直接替换原计划}——
这正是 V1 策略费用比问题二策略高出 24.4\% 的原因。

\textbf{（2）后续预报的价值主要体现在“触发调整”与“参与标定”。} 即使原始预报质量不佳，
以它为输入的双向调整机制（V2）仍能挽回 63\% 的预报损失；
而经过 MOS 标定并与 w1 融合后（V3、V4FR），后续预报才真正转化为经济收益。

\textbf{（3）当日已实现信息是主要的日内信息增量。} 以已实现的负载与光伏偏差按适度权重修正后续估计
（$\lambda_{pv}=0.1$、$\lambda_L=0.3$），可使总费用再降低 7.5 万元；
权重过大则会追逐噪声、增大调整摩擦，说明残差更新应当保守使用。

\textbf{（4）调整机制的收益来自费率结构的改善。} 0.5 倍电价的双向调整使模型在
“少购违约（$0.5c_t$）”“多购溢价（$1.5c_t$）”与“紧急购电（$5c_t$）”三个通道之间自动比价，
把问题二中被动承担的高成本缺口转化为主动的、低成本的计划调整。

综上，本文认为\textbf{应当引入其他时刻的预报，但应将其作为经标定与融合后的输入，
并配合当日已实现残差的适度修正与末端备用约束统一纳入滚动调整模型}，
而不是直接以最新预报取代原计划。
